% TOP_PROBLEM: {"id": 3151, "title": "Graham's conjecture on tree reconstruction", "category_id": 3, "category": {"id": 3, "name": "graph_theory", "display_name": "Graph Theory", "description": "Problems involving graphs, networks, and their properties.", "slug": "graph-theory", "order_index": 3, "created_at": "2026-07-31T15:26:25.671Z"}, "status": "open", "problem_number": "OPG-164", "set_id": 12}
% FIRST_POSTED: 2026-10-01
% ATTEMPT_STATUS: unresolved
% UnsolvedMath ID 3151; source catalogue code OPG-164
% !TeX program = lualatex
% Research attempt by GPT-6 Astra Ultra; no claim of a new complete solution.
\documentclass[11pt,a4paper]{article}
\usepackage[margin=25mm]{geometry}
\usepackage{fontspec}
\setmainfont{lmroman10-regular.otf}[BoldFont=lmroman10-bold.otf,
 ItalicFont=lmroman10-italic.otf,BoldItalicFont=lmroman10-bolditalic.otf]
\usepackage{amsmath,amssymb,mathtools}
\usepackage{unicode-math}
\setmathfont{latinmodern-math.otf}
\AtBeginDocument{\let\setminus\smallsetminus}
\usepackage{xurl,hyperref}
\hypersetup{colorlinks=true,urlcolor=blue}
\setcounter{secnumdepth}{0}
\setlength{\parindent}{0pt}
\setlength{\parskip}{5pt}
\setlength{\emergencystretch}{3em}
\begin{document}
\section*{Graham's conjecture on tree reconstruction}\label{3151}
{\small UnsolvedMath ID 3151 · OPG-164 · Graph Theory · OpenGarden}
\subsection{Definitions and mathematical statement}
For a finite simple graph $G$, its line graph $L(G)$ has one vertex
for each edge of $G$. Two distinct vertices of $L(G)$ are adjacent
precisely when the corresponding edges of $G$ share an endpoint.
Define iterates by $L^0(G)=G$ and $L^{j+1}(G)=L(L^j(G))$ for every
integer $j\geq0$. Each iterate is again a finite simple graph.
The line graph of an edgeless graph has no vertices, and subsequent
iterates stay empty.

Associate with $G$ its infinite sequence of orders
\[
 a(G)=(a_0(G),a_1(G),\ldots),\qquad a_j(G)=|V(L^j(G))|.
\]
A tree is a nonempty connected graph with no cycles. Graham's
reconstruction question is whether
\[
 \bigl[a_j(T)=a_j(T')\text{ for every }j\geq0\bigr]
 \quad\Longrightarrow\quad T\cong T'
\]
for every pair of finite trees. Isomorphism allows any renaming
of vertices and must preserve the edge relation.

Only the numbers of vertices are supplied, not the line graphs
themselves, their degree sequences, or maps between successive
iterates. The first two entries of an $n$-vertex tree are $n$ and
$n-1$, while later entries reflect more structure. The entire
infinite sequence is part of the statement: a collision among
finitely many initial entries alone does not disprove the
conjecture. A negative answer requires nonisomorphic trees with
agreement at every iteration.
\subsection{Short English statement}
Can a finite tree be reconstructed, up to isomorphism, from the vertex counts of all its iterated line graphs?
\subsection{Research attempt}
\paragraph{Scope and process.}
This is a GPT-6 Astra Ultra research attempt dated 1 October 2026, using
primary-source browsing and elementary mathematical checks. The canonical
definition above is retained. Results below are restricted results or
reductions, not a claim of a new solution to the entire catalogue question.
No formal proof assistant or external mathematical referee checked this text.


\paragraph{Literature and scope.}
Cooper, Kay and Swifton [R1] obtain a large family of distinguishable
trees using partition constructions. Weatherspoon and Zeilberger [R2]
report computational verification through order ten. Such finite-order
verification does not provide a uniform reconstruction theorem. The
attempt here derives the first structural terms, reconstructs the
extremal families from them, and gives two explicit trees with the same
first four terms but different fifth terms. This tests precisely how
much information a short-prefix approach discards.

\paragraph{The first terms and their exact meaning.}
Let $T$ have $n\ge2$ vertices and degrees $d_v$. The first two entries
are $a_0=n$ and $a_1=n-1$. Every edge of $L(T)$ corresponds to two
edges of $T$ meeting at a unique vertex, so
\[
 a_2=\sum_v\binom{d_v}{2}.
\]
An edge $uv$ of a simple graph becomes a vertex of its line graph of
degree $d_u+d_v-2$. Applying the same edge-count identity once more gives
\[
 a_3=\sum_{uv\in E(T)}\binom{d_u+d_v-2}{2}.
\]
Thus $a_2$ sees a degree moment, while $a_3$ also sees how endpoint
degrees correlate along edges. It is not legitimate to regard $a_2$
as determining the whole degree sequence or $a_3$ as specifying all
these individual correlations: each formula supplies only one sum.

\paragraph{A complete extremal reconstruction from $a_2$.}
Put $x_v=d_v-1\ge0$. The tree degree sum implies $\sum_vx_v=n-2$.
Since $\binom{d_v}{2}=x_v+\binom{x_v}{2}$,
\[
 a_2=n-2+\sum_v\binom{d_v-1}{2}\ge n-2.
\]
Equality holds exactly when every $d_v\le2$. A connected acyclic graph
with that property is a path, so paths are recognizable among all trees
by $a_0$ and $a_2$ alone. This includes $n=2$.

At the other extreme, $\sum_vx_v^2\le(\sum_vx_v)^2$ because all cross
products are nonnegative. Therefore
\[
 a_2=\tfrac12\sum_v(x_v^2+x_v)
 \le\tfrac12((n-2)^2+n-2)=\binom{n-1}{2}.
\]
For $n\ge3$, equality holds precisely when at most one $x_v$ is
positive. Consequently one vertex has degree $n-1$ and every other
vertex is a leaf: the tree is a star. The tree of order one is uniquely
recognized by $a_0=1$, and its later terms vanish. These boundary cases
ensure that the restricted reconstruction is valid at every order.

\paragraph{A checked collision beyond the degree moment.}
Consider the following eight-vertex trees, where $ij$ denotes an edge:
\[
 \begin{split}
 E(T_1)&=\{01,02,03,14,15,26,37\},\\
 E(T_2)&=\{01,02,03,14,15,26,47\}.
 \end{split}
\]
Both have degree multiset $(3,3,2,2,1,1,1,1)$. In the first, one
degree-three vertex has two degree-two neighbours and the other has
none; in the second, each degree-three vertex has exactly one
degree-two neighbour. This proves nonisomorphism without depending on
the vertex labels.

Direct line-graph construction gives
\[
 (a_0,a_1,a_2,a_3,a_4)(T_1)=(8,7,8,14,41),\qquad
 (a_0,a_1,a_2,a_3,a_4)(T_2)=(8,7,8,14,40).
\]
The first four terms therefore do not suffice, even at order eight.
The fifth term repairs this particular collision. The example is not
evidence of equality at every iteration, and calling it a counterexample
to Graham's conjecture would reverse the point of the computation.

\paragraph{Verification and limits of the attack.}
The script \texttt{scripts/check\_attempt\_batch02\_algebra\_2026\_10\_01.py}
constructs each line graph from its edge-incidence definition through
the fourth iterate. It checks the displayed sequences and independently
checks the formulas for $a_2,a_3$ on the two original trees. The test is
exact integer arithmetic and does not substitute a degree-based
approximation for the later line graphs.

A fixed finite prefix can separate a chosen finite list without supplying
an a priori stopping rule for all trees. Conversely, finding long
collisions says nothing about agreement of the entire infinite sequence
unless an invariant or recurrence proves persistence. The terms grow
rapidly for branching examples, so simply extending direct graph
construction is also a poor route to an unrestricted proof.

\paragraph{Final status and first remaining gap.}
\textbf{UNRESOLVED.} Paths and stars are reconstructed exactly; a short
prefix fails on a concrete nonisomorphic pair. What remains is a theorem
that some finite iterate distinguishes every unequal pair, or a pair
whose equality persists under all iterates. Neither the moment identities
nor the finite computation supplies that persistence argument.

\subsection{Sources}
\begin{itemize}
\item[S1] ULAM.AI and the UnsolvedMath Contributors, supplied problems.json, record 3151 (OPG-164), OpenGarden collection. Catalogue identity and source transcription; CC BY 4.0.
\url{https://huggingface.co/datasets/ulamai/UnsolvedMath}.
\item[S2] Open Problem Garden, Graham's conjecture on tree reconstruction. Original problem and discussion; the mathematical formulation below is expanded from the supplied source record.
\url{http://www.openproblemgarden.org/op/grahams_conjecture_on_tree_reconstruction}.

\item[R1] Joshua Cooper, Bill Kay and Anton Swifton,
\emph{Graham's Tree Reconstruction Conjecture and a Waring-Type Problem
on Partitions}, arXiv version revised 2017. Primary abstract checked
1 October 2026.
\url{https://arxiv.org/abs/1109.0522}.
\item[R2] Kaylee Weatherspoon and Doron Zeilberger,
\emph{An Experimental Note on Graham's Tree Reconstruction Conjecture},
June 2025. Primary author-hosted computational note, checked 1 October 2026.
\url{https://sites.math.rutgers.edu/~zeilberg/mamarim/mamarimPDF/treeline.pdf}.

\end{itemize}
\end{document}
